% HiddenPower — итерация 6, накопительная версия V06.
% Глава: «Уравнения в частных производных и дифференцируемые полевые модели».
% Источники: J. K. Hunter; C. Clason.

\chapter{Уравнения в частных производных и дифференцируемые полевые модели}
\label{chap:partial-differential-equations}

В главах~\ref{chap:dynamics-differentiable-models} и~\ref{chap:stochastic-differential-equations} состояние зависело от времени. Для конструкции этого часто недостаточно: температура, смещение, напряжение, плотность материала и пористость распределены по пространству. Вместо конечномерного вектора возникает поле
\[
u:D\times I\to\R^m,
\qquad
(x,t)\longmapsto u(x,t),
\]
где $D\subset\R^d$ --- пространственная область, $I$ --- интервал времени, а $m$ задаёт число компонент поля.

Основная вычислительная цепочка главы имеет вид
\[
\begin{aligned}
\text{геометрия и материал}
&\longrightarrow
\text{область, коэффициенты и данные},\\
\text{PDE и граничные условия}
&\longrightarrow
\text{дискретный оператор},\\
\text{дискретный оператор}
&\longrightarrow
\text{поле отклика}
\longrightarrow
\text{целевой функционал и градиент}.
\end{aligned}
\]
Эта цепочка отделяет физическую постановку от численного решателя. В первых версиях главы уточняется, какое поле неизвестно, где оно определено и какие данные делают задачу осмысленной. Затем PDE переводится в конечномерную систему, пригодную для вычислений и дифференцирования по параметрам конструкции.

\section{Поля, уравнения в частных производных и постановка задачи}
\label{sec:pde-fields-and-problem-statement}

\subsection{Пространственное поле как неизвестная функция}

Пусть $D\subset\R^d$ --- открытая область. Чтобы не смешивать её с пространством случайных исходов $\Omega$ из главы~\ref{chap:stochastic-differential-equations}, пространственную область во всей главе обозначаем буквой $D$. Точка $x=(x_1,\ldots,x_d)\in D$ задаёт положение в конструкции.

Скалярное поле имеет вид
\[
u:D\to\R.
\]
Так можно описывать температуру, концентрацию, плотность материала или потенциал. Векторное поле
\[
u:D\to\R^m
\]
описывает, например, перемещение каждой точки конструкции. Для нестационарной модели добавляется время:
\[
u:D\times(0,T)\to\R^m.
\]

\begin{definition}
Уравнением в частных производных называется соотношение
\begin{equation}
\label{eq:general-pde-v01}
F\!\left(
 x,t,u,\partial u,\partial^2u,\ldots,\partial^k u
\right)=0,
\end{equation}
связывающее неизвестное поле $u$ с его частными производными по нескольким независимым переменным.
\end{definition}

Отличие от ОДУ состоит не только в числе переменных. Для ОДУ начальное состояние обычно задаётся в одной точке времени. Для PDE неизвестна функция на целой области, а данные задаются на границе пространства, в начальный момент или одновременно на нескольких частях пространственно-временной границы.

\begin{application}[поля в генеративном проектировании]
Пусть параметр $\theta$ задаёт форму области $D(\theta)$ и распределение материала $a(x;\theta)$. Поле $u(x;\theta)$ может быть температурой или смещением, найденным из PDE. Генеративная модель предлагает $\theta$, но качество кандидата определяется не самим кодом, а вычисленным полем отклика и функционалом
\begin{equation}
\label{eq:pde-design-objective-preview-v01}
J(\theta)=\mathcal J\bigl(u(\cdot;\theta),\theta\bigr).
\end{equation}
Поэтому физический решатель является отображением от описания конструкции к её распределённому отклику.
\end{application}

\subsection{Порядок, линейность и старшая часть}

Порядок PDE равен наибольшему порядку производной, входящей в уравнение. Уравнение
\[
-\nabla\cdot\bigl(a(x)\nabla u(x)\bigr)=f(x)
\]
имеет второй порядок, поскольку после раскрытия дивергенции появляются вторые производные $u$. Уравнение переноса
\[
u_t+b(x,t)^{\trans}\nabla u=0
\]
имеет первый порядок.

Линейное скалярное уравнение второго порядка можно записать как
\begin{equation}
\label{eq:linear-second-order-pde-v01}
\mathcal Lu
=
-\sum_{i,j=1}^{d}
\partial_i\!\left(a_{ij}(x,t)\partial_j u\right)
+\sum_{i=1}^{d}b_i(x,t)\partial_i u
+c(x,t)u
=f.
\end{equation}
Коэффициенты и правая часть заданы, а $u$ входит линейно. Если коэффициент зависит от самого решения, например
\[
-\nabla\cdot\bigl(a(u)\nabla u\bigr)=f,
\]
уравнение становится нелинейным.

Матрица
\begin{equation}
\label{eq:principal-matrix-v01}
A(x,t)=\bigl(a_{ij}(x,t)\bigr)_{i,j=1}^{d}
\end{equation}
задаёт старшую часть оператора. Именно она определяет основные свойства PDE: направления распространения, сглаживание, устойчивость и тип допустимых граничных данных. Члены с первыми производными и без производных важны, но не меняют порядок уравнения.

Если $A$ симметрична, то квадратичная форма
\[
\xi^{\trans}A(x,t)\xi
\]
показывает силу действия оператора вдоль направления $\xi$. Эта связь возвращает нас к спектральному анализу из главы~\ref{chap:linear-algebra}: собственные векторы задают главные направления, а собственные значения --- относительные масштабы диффузии или проводимости.

\subsection{Область, граница и данные задачи}

Само дифференциальное уравнение редко определяет единственное поле. Нужно указать область $D$, правую часть, коэффициенты и данные на границе $\partial D$. Для нестационарной задачи также задаются данные при $t=0$.

Условие Дирихле фиксирует значение поля:
\begin{equation}
\label{eq:dirichlet-condition-v01}
u=g_D
\qquad
\text{на }\Gamma_D\subseteq\partial D.
\end{equation}
Для температуры это заданная температура на поверхности, для перемещения --- закрепление точки или части границы.

Условие Неймана фиксирует нормальный поток. Для оператора с матрицей $A$ естественная запись имеет вид
\begin{equation}
\label{eq:neumann-condition-v01}
 n^{\trans}A\nabla u=g_N
\qquad
\text{на }\Gamma_N\subseteq\partial D,
\end{equation}
где $n$ --- внешняя единичная нормаль. При $A=I$ получается обычная нормальная производная $\partial_nu=n^{\trans}\nabla u$.

Условие Робина связывает значение и поток:
\begin{equation}
\label{eq:robin-condition-v01}
\alpha u+\beta\,n^{\trans}A\nabla u=g_R
\qquad
\text{на }\Gamma_R.
\end{equation}
Так моделируется, например, теплообмен с окружающей средой. Коэффициенты $\alpha$ и $\beta$ имеют физические размерности; их нельзя выбирать как безразмерные веса только ради удобства оптимизации.

Для уравнения первого порядка по времени обычно требуется одно начальное поле
\begin{equation}
\label{eq:first-order-initial-condition-v01}
u(x,0)=u_0(x).
\end{equation}
Для уравнения второго порядка по времени, как в волновой задаче, задаются положение и скорость:
\begin{equation}
\label{eq:second-order-initial-condition-v01}
u(x,0)=u_0(x),
\qquad
u_t(x,0)=v_0(x).
\end{equation}
Число начальных условий согласуется с порядком производной по времени.

\begin{caution}[граница является частью модели]
Одинаковое PDE с разными граничными условиями описывает разные физические системы. Закреплённая граница, свободная граница и заданная нагрузка не являются вариантами оформления одной задачи. Они меняют множество допустимых решений и после дискретизации меняют матрицу и правую часть системы.
\end{caution}

\subsection{Эллиптические, параболические и гиперболические задачи}

Три основных типа удобно различать по характеру старшей части и по роли времени. Эллиптическая задача описывает стационарное равновесие:
\begin{equation}
\label{eq:elliptic-model-v01}
-\nabla\cdot\bigl(A(x)\nabla u\bigr)=f
\qquad
\text{в }D.
\end{equation}
Если существует $\lambda>0$ такое, что
\begin{equation}
\label{eq:uniform-ellipticity-v01}
\xi^{\trans}A(x)\xi
\geq
\lambda\norm{\xi}^2
\qquad
\text{для всех }x\in D,\ \xi\in\R^d,
\end{equation}
оператор называется равномерно эллиптическим. Это условие запрещает направления, в которых старшая часть полностью исчезает.

Параболическая задача описывает необратимое сглаживание во времени. Канонический пример --- уравнение теплопроводности
\begin{equation}
\label{eq:parabolic-model-v01}
u_t-\nabla\cdot\bigl(A(x,t)\nabla u\bigr)=f.
\end{equation}
Здесь время входит первой производной, а пространственная часть остаётся эллиптической.

Гиперболическая задача описывает распространение волн с конечной скоростью. Канонический пример имеет вид
\begin{equation}
\label{eq:hyperbolic-model-v01}
u_{tt}-c^2\Delta u=f.
\end{equation}
В отличие от теплопроводности, волновое уравнение не сглаживает начальные данные тем же способом: локальное возмущение распространяется по характеристическим направлениям.

\begin{problem}[законченный расчёт классификации оператора]
Рассмотрим стационарное уравнение
\begin{equation}
\label{eq:classification-problem-v01}
-2u_{xx}-2u_{xy}-2u_{yy}=f(x,y).
\end{equation}
Определим тип его старшей части.
\end{problem}

Смешанная производная в квадратичной форме учитывается дважды, поэтому матрица старшей части равна
\[
A=
\begin{pmatrix}
2&1\\
1&2
\end{pmatrix},
\qquad
-\nabla\cdot(A\nabla u)
=-2u_{xx}-2u_{xy}-2u_{yy}.
\]
Вычисляем собственные значения:
\begin{align}
\det(A-\lambda I)
&=
\det
\begin{pmatrix}
2-\lambda&1\\
1&2-\lambda
\end{pmatrix}
\\
&=(2-\lambda)^2-1
\\
&=(\lambda-1)(\lambda-3).
\end{align}
Следовательно,
\[
\lambda_1=1,
\qquad
\lambda_2=3.
\]
Оба значения положительны. Для любого $\xi\in\R^2$
\[
\xi^{\trans}A\xi
\geq
\lambda_{\min}(A)\norm{\xi}^2
=
\norm{\xi}^2.
\]
Условие~\eqref{eq:uniform-ellipticity-v01} выполнено с константой $\lambda=1$, поэтому уравнение~\eqref{eq:classification-problem-v01} является равномерно эллиптическим. Расчёт показывает практический порядок проверки: выписать симметричную матрицу старшей части, найти её спектр и проверить знак наименьшего собственного значения.

\subsection{Корректность и совместимость условий}

Постановка считается корректной, если решение существует, единственно в выбранном классе и непрерывно зависит от данных. Последнее свойство означает, что малая ошибка в нагрузке, коэффициентах или граничных значениях не должна вызывать неограниченное изменение решения. Типичная оценка имеет форму
\begin{equation}
\label{eq:pde-stability-estimate-v01}
\norm{u-\widetilde u}_{X}
\leq
C\left(
\norm{f-\widetilde f}_{Y}
+
\norm{g-\widetilde g}_{Z}
\right),
\end{equation}
где пространства $X,Y,Z$ и константа $C$ зависят от конкретного PDE и граничной задачи.

Существование без единственности недостаточно для проектирования: один набор параметров тогда соответствует нескольким откликам. Единственность без устойчивости также недостаточна: вычислительная ошибка или малая производственная вариация может полностью изменить результат.

Данные должны быть совместимы. Для задачи Неймана
\[
-\Delta u=f
\quad\text{в }D,
\qquad
\partial_nu=g
\quad\text{на }\partial D
\]
интегрирование и теорема о дивергенции дают
\[
\int_D f\,\dd x
=
-\int_D\Delta u\,\dd x
=
-\int_{\partial D}g\,\dd S.
\]
Поэтому необходимое условие имеет вид
\begin{equation}
\label{eq:neumann-compatibility-v01}
\int_D f\,\dd x
+
\int_{\partial D}g\,\dd S
=0.
\end{equation}
Даже при выполнении этого условия решение определяется только с точностью до добавления константы, поскольку градиент и лапласиан константы равны нулю. Для единственности нужно дополнительно зафиксировать, например, среднее значение $u$.

Для нестационарной задачи с условием Дирихле требуется согласование в начальный момент:
\[
u_0(x)=g_D(x,0)
\qquad
\text{на }\Gamma_D.
\]
Если геометрия имеет углы, разрывы коэффициентов или несовместимые данные, гладкого классического решения может не существовать. Это не всегда означает, что физическая задача бессмысленна: позже будет введено слабое решение, допускающее меньшую регулярность.

\subsection{От физической постановки к вычислимой модели}

Для генеративного проектирования полезно фиксировать постановку в следующем порядке:
\[
\begin{aligned}
\theta
&\longmapsto
D(\theta),\ A(x;\theta),\ f(x;\theta),\ g(x;\theta),\\
\bigl(D,A,f,g\bigr)
&\longmapsto
u(\cdot;\theta),\\
u(\cdot;\theta)
&\longmapsto
J\bigl(u,\theta\bigr).
\end{aligned}
\]
Параметр $\theta$ может кодировать толщину, ориентацию материала, пористость или форму отверстий. В следующих версиях поле будет заменено вектором $U$, а оператор --- матрицей $K(\theta)$.

\begin{mlconnection}[что именно получает модель машинного обучения]
Входом служат геометрия, поле материала и нагрузка; выходом --- конкретное поле $u$ или функционал $J$. Малая невязка PDE не гарантирует правильных граничных условий, точности сетки и корректности исходной физической модели.
\end{mlconnection}

Градиент и дивергенция опираются на главу~\ref{chap:multivariable-calculus}, спектр старшей части --- на главу~\ref{chap:linear-algebra}. Для полей на поверхностях нужны конструкции главы~\ref{chap:manifolds-differential-geometry}, а временные PDE продолжают главу~\ref{chap:dynamics-differentiable-models}.

Ограничения фиксируют до оптимизации. Континуальная модель ненадёжна на масштабе дискретной структуры материала; трещины, контакт и изменение топологии создают негладкости, а плохая сетка ухудшает обусловленность. Даже точное решение PDE бесполезно при неверных коэффициентах, нагрузках или граничных условиях.


\section{Перенос и законы сохранения}
\label{sec:pde-transport-and-conservation}

\subsection{От интегрального баланса к локальному уравнению}

Пусть $q(x,t)$ --- плотность некоторой величины: массы, энергии, концентрации или объёмной доли материала. Чтобы модель действительно выражала сохранение, недостаточно написать произвольное PDE. Сначала фиксируется баланс на произвольной подобласти $V\Subset D$.

Обозначим через $J(x,t)\in\R^d$ поток величины $q$, а через $s(x,t)$ --- объёмную скорость её создания. Внешняя единичная нормаль к $\partial V$ обозначается $n$. Интегральный баланс имеет вид
\begin{equation}
\label{eq:integral-balance-v02}
\frac{\dd}{\dd t}
\int_V q(x,t)\,\dd x
=
-\int_{\partial V}J(x,t)^{\trans}n\,\dd S
+
\int_V s(x,t)\,\dd x.
\end{equation}
Знак минус перед потоком означает, что положительный внешний поток уменьшает количество величины внутри $V$.

Если $q$ и $J$ достаточно регулярны, теорема о дивергенции даёт
\[
\int_{\partial V}J^{\trans}n\,\dd S
=
\int_V\nabla\cdot J\,\dd x.
\]
Подставляя это равенство в~\eqref{eq:integral-balance-v02}, получаем
\[
\int_V
\left(
q_t+\nabla\cdot J-s
\right)\dd x
=0.
\]
Так как подобласть $V$ выбирается произвольно, локальный закон сохранения принимает вид
\begin{equation}
\label{eq:local-conservation-law-v02}
q_t+\nabla\cdot J=s
\qquad
\text{в }D\times(0,T).
\end{equation}

Уравнение~\eqref{eq:local-conservation-law-v02} показывает, какие объекты нужно определить до вычислений. Неизвестной является плотность $q$, физическая модель задаёт поток $J$, а источники и стоки входят через $s$. Если поток выбран неверно, формальное выполнение PDE не исправляет физическую постановку.

\begin{application}[плотность материала в проектируемой области]
Пусть $q=\rho(x,t)\in[0,1]$ кодирует локальную долю материала. Интеграл
\[
M(t)=\int_D\rho(x,t)\,\dd x
\]
пропорционален объёму материала. При $s=0$ и условии $J^{\trans}n=0$ на $\partial D$ из~\eqref{eq:integral-balance-v02} следует $M'(t)=0$. Поэтому сохранение общего количества материала проверяется интегралом, а не только малой поточечной невязкой PDE.
\end{application}

\subsection{Адвективный поток и две формы уравнения}

Пусть задано поле скорости $\beta(x,t)\in\R^d$. Если величина переносится вместе с этим полем, естественный поток равен
\begin{equation}
\label{eq:advective-flux-v02}
J=\beta q.
\end{equation}
Тогда закон сохранения~\eqref{eq:local-conservation-law-v02} становится
\begin{equation}
\label{eq:conservative-transport-v02}
q_t+\nabla\cdot(\beta q)=s.
\end{equation}
По правилу произведения
\begin{equation}
\label{eq:divergence-product-v02}
\nabla\cdot(\beta q)
=
\beta^{\trans}\nabla q
+
(\nabla\cdot\beta)q.
\end{equation}
Следовательно,
\begin{equation}
\label{eq:expanded-conservative-transport-v02}
q_t
+
\beta^{\trans}\nabla q
+
(\nabla\cdot\beta)q
=s.
\end{equation}

Нельзя без проверки заменять~\eqref{eq:conservative-transport-v02} на
\begin{equation}
\label{eq:nonconservative-transport-v02}
u_t+\beta^{\trans}\nabla u=r.
\end{equation}
Уравнение~\eqref{eq:nonconservative-transport-v02} описывает изменение значения, наблюдаемого вдоль движущейся траектории. Уравнение~\eqref{eq:conservative-transport-v02} описывает плотность, для которой важен также локальный объём потока.

Если $\nabla\cdot\beta=0$, поле скорости несжимаемо и две формы совпадают при $u=q$ и $r=s$. Если $\nabla\cdot\beta\neq0$, расширение или сжатие потока меняет плотность даже при отсутствии источника. Именно член $(\nabla\cdot\beta)q$ обеспечивает корректный баланс.

\begin{caution}[перенос значения не равен переносу плотности]
Для цвета, температуры-маркера или иной пассивной метки часто естественно уравнение~\eqref{eq:nonconservative-transport-v02}. Для массы и объёмной доли материала нужна консервативная форма~\eqref{eq:conservative-transport-v02}. Выбор между ними определяется смыслом поля, а не удобством реализации.
\end{caution}

\subsection{Характеристики и связь с обыкновенными уравнениями}

Рассмотрим уравнение~\eqref{eq:nonconservative-transport-v02}. Характеристическая траектория $X(t)$ определяется ОДУ
\begin{equation}
\label{eq:characteristic-ode-v02}
\dot X(t)=\beta(X(t),t),
\qquad
X(0)=\xi.
\end{equation}
Если поле $\beta$ удовлетворяет условиям существования и единственности из главы~\ref{chap:dynamics-differentiable-models}, то каждой начальной точке $\xi$ соответствует единственная траектория.

По правилу цепочки
\begin{align}
\frac{\dd}{\dd t}u(X(t),t)
&=
 u_t(X(t),t)
+
\dot X(t)^{\trans}\nabla u(X(t),t)
\\
&=
 u_t(X(t),t)
+
\beta(X(t),t)^{\trans}\nabla u(X(t),t).
\end{align}
Поэтому на характеристике уравнение~\eqref{eq:nonconservative-transport-v02} сводится к ОДУ
\begin{equation}
\label{eq:value-along-characteristic-v02}
\frac{\dd}{\dd t}u(X(t),t)=r(X(t),t).
\end{equation}
При $r=0$ значение $u$ постоянно вдоль траектории. PDE не создаёт новый профиль, а переносит начальные значения по потоку.

Для постоянной скорости $\beta(x,t)=c\in\R^d$ характеристика имеет вид
\[
X(t)=\xi+ct.
\]
Из $x=X(t)$ получаем $\xi=x-ct$. Если $u(x,0)=u_0(x)$ и $r=0$, решение равно
\begin{equation}
\label{eq:constant-transport-solution-v02}
u(x,t)=u_0(x-ct).
\end{equation}
Это формула сдвига: форма профиля сохраняется, а его положение меняется со скоростью $c$.

\begin{problem}[законченный расчёт переноса профиля]
Дано одномерное уравнение
\begin{equation}
\label{eq:transport-calculation-problem-v02}
u_t+2u_x=0,
\qquad
u(x,0)=\exp\bigl(-(x-1)^2\bigr),
\qquad
x\in\R.
\end{equation}
Найдём решение, значение $u(2,0.75)$ и проверим сохранение интеграла профиля.
\end{problem}

Скорость равна $c=2$. По формуле~\eqref{eq:constant-transport-solution-v02}
\begin{equation}
\label{eq:transport-calculation-solution-v02}
u(x,t)
=
\exp\bigl(-(x-2t-1)^2\bigr).
\end{equation}
При $x=2$ и $t=0.75$ сначала вычисляем аргумент начального профиля:
\[
x-ct=2-2\cdot0.75=0.5.
\]
Затем подставляем его в $u_0$:
\begin{align}
u(2,0.75)
&=
\exp\bigl(-(0.5-1)^2\bigr)
\\
&=
\exp(-0.25)
\\
&\approx0.7788.
\end{align}
Центр начального профиля находился в точке $x=1$. За время $0.75$ он сместился на $2\cdot0.75=1.5$, поэтому новый центр равен $2.5$.

Для проверки сохранения полного количества делаем замену переменной $y=x-2t$:
\begin{align}
\int_{\R}u(x,t)\,\dd x
&=
\int_{\R}u_0(x-2t)\,\dd x
\\
&=
\int_{\R}u_0(y)\,\dd y
\\
&=
\int_{\R}\exp\bigl(-(y-1)^2\bigr)\,\dd y
\\
&=
\sqrt{\pi}.
\end{align}
Итак, профиль сдвинулся без изменения формы и интеграла. Расчёт одновременно проверяет решение PDE и его консервативный смысл при постоянном поле скорости.

\subsection{Входящая и выходящая части границы}

Для уравнения первого порядка направление потока определяет, где можно задавать граничные данные. Пусть $n(x)$ --- внешняя нормаль к $\partial D$. Вводятся множества
\begin{align}
\Gamma_-
&=
\left\{
 x\in\partial D:
 \beta(x)^{\trans}n(x)<0
\right\},
\label{eq:inflow-boundary-v02}
\\
\Gamma_+
&=
\left\{
 x\in\partial D:
 \beta(x)^{\trans}n(x)>0
\right\}.
\label{eq:outflow-boundary-v02}
\end{align}
На $\Gamma_-$ характеристики входят в область, поэтому значение решения должно быть задано там. На $\Gamma_+$ характеристики выходят, и значение определяется внутренним решением. Одновременное независимое задание данных на обеих частях обычно переопределяет задачу.

Clason рассматривает стационарное уравнение переноса с реакцией
\begin{equation}
\label{eq:stationary-advection-reaction-v02}
\beta^{\trans}\nabla u+\mu u=f
\qquad
\text{в }D,
\end{equation}
дополненное условием на входящей границе. Коэффициент $\mu$ задаёт локальное создание или уничтожение переносимой величины, а $f$ --- внешний источник.

Естественное пространство должно контролировать не весь градиент, а производную вдоль потока:
\begin{equation}
\label{eq:transport-graph-space-v02}
W
=
\left\{
 v\in L^2(D):
 \beta^{\trans}\nabla v\in L^2(D)
\right\}.
\end{equation}
Соответствующая норма имеет вид
\begin{equation}
\label{eq:transport-graph-norm-v02}
\norm{v}_{W}^2
=
\norm{v}_{L^2(D)}^2
+
\norm{\beta^{\trans}\nabla v}_{L^2(D)}^2.
\end{equation}

Для устойчивости стационарной задачи достаточно условия
\begin{equation}
\label{eq:advection-reaction-stability-v02}
\mu(x)-\frac12\nabla\cdot\beta(x)
\geq
\mu_0>0.
\end{equation}
Оно показывает, что знак реакции нужно оценивать совместно со сжатием или расширением поля скорости. Положительное $\mu$ само по себе не гарантирует ту же оценку при большом положительном $\nabla\cdot\beta$.

\subsection{Перенос как преобразование конструкции}

В генеративном проектировании поле скорости может использоваться не как физическая скорость, а как вычислительное правило деформации. Пусть $\rho_0(x)$ задаёт исходное распределение материала, а $\beta(x,t;\theta)$ зависит от параметров модели. Консервативная динамика
\begin{equation}
\label{eq:design-density-transport-v02}
\rho_t
+
\nabla\cdot\bigl(\beta(x,t;\theta)\rho\bigr)
=0
\end{equation}
перемещает материал по области. Если $\beta^{\trans}n=0$ на $\partial D$, то
\[
\frac{\dd}{\dd t}
\int_D\rho(x,t)\,\dd x
=0.
\]
Такой слой может изменять расположение материала, не меняя его суммарное количество.

Вычислительная цепочка здесь конкретна:
\[
\theta
\longmapsto
\beta(\cdot,\cdot;\theta)
\longmapsto
\rho(\cdot,T;\theta)
\longmapsto
\text{геометрия}
\longmapsto
J(\theta).
\]
Параметры $\theta$ могут задавать поле деформации, а PDE превращает это поле в пространственно согласованное изменение конструкции. В V06 производная $\dd J/\dd\theta$ будет вычисляться через дискретный решатель.

\begin{mlconnection}[проверяемая роль модели]
Модель машинного обучения может предсказывать $\beta$, источник $s$ или конечное поле. Физическая проверка состоит в балансе~\eqref{eq:integral-balance-v02}, корректном условии на $\Gamma_-$ и соответствии потока смыслу переменной. Само наличие PDE-невязки в функции потерь не гарантирует сохранения после дискретизации, интерполяции и отсечения значений.
\end{mlconnection}

Чистый перенос имеет ограниченную выразительность. При гладком однозначном потоке он деформирует уже существующую структуру, но не обязан создавать или удалять компоненты связности. Если генеративная задача требует рождения отверстий, разрыва перемычек или изменения числа компонентов, одного уравнения переноса может быть недостаточно.

\subsection{Численные ограничения и границы применимости}

Транспортное уравнение направлено вдоль $\beta$, поэтому дискретизация должна учитывать входящую и выходящую стороны элементов. Clason использует разрывные конечномерные пространства как естественную основу для уравнений первого порядка: решение допускается разрывным между элементами, а обмен информацией задаётся через потоки на гранях.

Гладкость решения определяется не только коэффициентами. Если $u_0$ имеет скачок, формула~\eqref{eq:constant-transport-solution-v02} переносит этот скачок без сглаживания. Классическая производная на линии разрыва не существует, и нужна слабая постановка. Численная схема может либо создавать колебания около скачка, либо чрезмерно размазывать его.

Для характеристического описания поле $\beta$ должно обеспечивать существование однозначного потока. При недостаточной регулярности траектории могут перестать быть единственными, и формула переноса начальных данных теряет однозначность.

Отдельно проверяются размерности и масштабирование. Если координаты измеряются в миллиметрах, а скорость --- в метрах в секунду, коэффициенты нельзя подставлять без преобразования единиц. При использовании нормированных координат нужно вернуть физический масштаб перед интерпретацией потока и времени.

Проекция $\rho\mapsto\min(1,\max(0,\rho))$ обеспечивает диапазон $[0,1]$, но в общем случае меняет интеграл материала. Поэтому ограничение диапазона и сохранение массы являются разными требованиями и должны проверяться отдельно.

Раздел продолжает правило цепочки из главы~\ref{chap:multivariable-calculus}, потоки ОДУ из главы~\ref{chap:dynamics-differentiable-models} и дивергентные операторы из главы~\ref{chap:manifolds-differential-geometry}. В следующих версиях локальный баланс будет дополнен диффузией, слабой постановкой и конечномерной аппроксимацией.

\section{Уравнения Лапласа и Пуассона}
\label{sec:laplace-poisson-equations}

\subsection{Равновесное поле и эллиптический оператор}

Уравнение Лапласа описывает поле без внутренних источников:
\begin{equation}
\label{eq:laplace-equation-v03}
-\Delta u=0
\qquad
\text{в }D.
\end{equation}
Знак перед лапласианом выбран так, чтобы оператор $-\Delta$ был положительным в энергетической постановке. Эквивалентная запись $\Delta u=0$ используется столь же часто. Решение уравнения Лапласа называется гармонической функцией.

Если в области действует распределённый источник $f$, возникает уравнение Пуассона
\begin{equation}
\label{eq:poisson-equation-v03}
-\Delta u=f
\qquad
\text{в }D.
\end{equation}
Для стационарной теплопроводности $u$ является температурой, а $f$ описывает внутреннее тепловыделение после согласования физических коэффициентов. В электростатике неизвестным может быть потенциал, а правая часть связана с плотностью заряда.

Неоднородный или анизотропный материал приводит к оператору
\begin{equation}
\label{eq:anisotropic-poisson-v03}
-\nabla\cdot\bigl(A(x)\nabla u(x)\bigr)=f(x),
\end{equation}
где $A(x)$ --- симметричная матрица проводимости. При $A(x)=a(x)I$ получаем скалярный коэффициент $a(x)$. Условие
\begin{equation}
\label{eq:anisotropic-ellipticity-v03}
\xi^{\trans}A(x)\xi
\geq
\lambda\norm{\xi}^{2},
\qquad
\lambda>0,
\end{equation}
исключает направления с нулевой проводимостью и обеспечивает контроль градиента решения.

Уравнения~\eqref{eq:laplace-equation-v03}--\eqref{eq:anisotropic-poisson-v03} описывают равновесие. В них нет производной по времени: поле уже достигло стационарного состояния. Это не означает отсутствия динамики в физической системе. Стационарное уравнение часто получается как предел нестационарного процесса, когда временная производная стала равна нулю.

\begin{connection}[связь с предыдущими главами]
Лапласиан объединяет операции градиента и дивергенции из глав~\ref{chap:multivariable-calculus} и~\ref{chap:manifolds-differential-geometry}:
\[
\Delta u=\nabla\cdot\nabla u.
\]
Матрица $A$ действует на градиент как линейный оператор, поэтому её спектр и обусловленность относятся к главе~\ref{chap:linear-algebra}. В отличие от траектории ОДУ из главы~\ref{chap:dynamics-differentiable-models}, решение здесь определяется одновременно уравнением во всей области и данными на её границе.
\end{connection}

\subsection{Граничные условия и формула Грина}

Рассмотрим оператор~\eqref{eq:anisotropic-poisson-v03}. Для достаточно гладких $u$ и $v$ интегрирование по частям даёт первую формулу Грина
\begin{equation}
\label{eq:green-identity-anisotropic-v03}
\int_D
v\bigl[-\nabla\cdot(A\nabla u)\bigr] \,\dd x
=
\int_D
(\nabla v)^{\trans}A\nabla u\,\dd x
-
\int_{\partial D}
v\,n^{\trans}A\nabla u\,\dd S.
\end{equation}
Она показывает, как производная второго порядка переносится с $u$ на тестовую функцию $v$ и почему поток появляется на границе.

Для условия Дирихле
\[
u=g_D
\qquad
\text{на }\partial D
\]
значение поля фиксируется непосредственно. При однородном условии $g_D=0$ допустимые вариации $v$ также обращаются в нуль на границе, поэтому граничный интеграл в~\eqref{eq:green-identity-anisotropic-v03} исчезает.

При условии Неймана
\[
n^{\trans}A\nabla u=g_N
\qquad
\text{на }\partial D
\]
фиксируется производная по внешней нормали с учётом коэффициента материала. Интегрирование уравнения по области даёт необходимое условие совместимости
\begin{equation}
\label{eq:neumann-compatibility-v03}
\int_D f\,\dd x
+
\int_{\partial D}g_N\,\dd S
=0.
\end{equation}
Знак следует из принятой записи $-\nabla\cdot(A\nabla u)=f$. Если физический выходящий поток определён как $q_n=-n^{\trans}A\nabla u$, формула записывается в эквивалентном виде
\[
\int_{\partial D}q_n\,\dd S
=
\int_D f\,\dd x.
\]

Чистая задача Неймана не определяет абсолютный уровень поля. Если $u$ является решением, то $u+C$ имеет тот же градиент и тот же поток. Единственность восстанавливают, например, нормировкой
\begin{equation}
\label{eq:neumann-zero-mean-v03}
\int_D u\,\dd x=0.
\end{equation}
В вычислительной системе отсутствие такой нормировки проявляется как нулевое собственное значение матрицы.

\subsection{Принцип максимума и единственность}

\begin{theorem}[принцип максимума]
Пусть $D$ --- ограниченная связная область, а функция $u\in C^2(D)\cap C(\overline D)$ гармонична. Тогда
\begin{equation}
\label{eq:maximum-principle-v03}
\min_{\partial D}u
\leq
u(x)
\leq
\max_{\partial D}u
\qquad
\text{для всех }x\in D.
\end{equation}
Если внутренний максимум или минимум достигается не постоянной гармонической функцией, возникает противоречие.
\end{theorem}

Практический смысл прост: стационарное поле без источников не может создать в глубине области температуру выше всех заданных граничных температур или ниже их минимума. Экстремальные значения контролируются границей.

Принцип максимума сразу даёт единственность задачи Дирихле. Пусть $u_1$ и $u_2$ решают одну задачу Пуассона с одинаковыми граничными данными. Разность
\[
w=u_1-u_2
\]
удовлетворяет
\[
\Delta w=0
\quad\text{в }D,
\qquad
w=0
\quad\text{на }\partial D.
\]
По~\eqref{eq:maximum-principle-v03} одновременно $w\leq0$ и $w\geq0$, следовательно, $w=0$ и $u_1=u_2$.

\begin{caution}[источник меняет вывод]
Принцип~\eqref{eq:maximum-principle-v03} нельзя без проверки переносить на уравнение Пуассона с ненулевой правой частью. Положительное внутреннее тепловыделение способно создать максимум температуры внутри области. Знак $f$, знак оператора и тип граничных условий должны учитываться совместно.
\end{caution}

Принцип максимума полезен и как численная диагностика. Если дискретное решение гармонической задачи с граничными значениями от $0$ до $1$ даёт внутреннее значение $1.4$, это не новый физический эффект, а признак ошибки постановки, сетки или схемы. Однако не всякая дискретизация автоматически сохраняет непрерывный принцип максимума.

\subsection{Энергетическая постановка}

Для симметричной положительно определённой матрицы $A$ задачу Дирихле можно связать с функционалом
\begin{equation}
\label{eq:poisson-energy-v03}
\mathcal E[v]
=
\frac12
\int_D
(\nabla v)^{\trans}A\nabla v\,\dd x
-
\int_D f v\,\dd x.
\end{equation}
Допустимые функции удовлетворяют заданному условию Дирихле. Пусть $u$ минимизирует $\mathcal E$, а вариация $\varphi$ равна нулю на границе. Для $v=u+\varepsilon\varphi$ имеем
\begin{align}
\left.
\frac{\dd}{\dd\varepsilon}
\mathcal E[u+\varepsilon\varphi]
\right|_{\varepsilon=0}
&=
\int_D
(\nabla\varphi)^{\trans}A\nabla u\,\dd x
-
\int_D f\varphi\,\dd x.
\end{align}
В точке минимума производная равна нулю:
\begin{equation}
\label{eq:poisson-variational-preview-v03}
\int_D
(\nabla\varphi)^{\trans}A\nabla u\,\dd x
=
\int_D f\varphi\,\dd x.
\end{equation}
После применения формулы Грина и учёта $\varphi|_{\partial D}=0$ получаем исходное PDE~\eqref{eq:anisotropic-poisson-v03}.

Эта запись важна по двум причинам. Во-первых, в ней достаточно первых производных $u$, тогда как классическая форма содержит вторые. Во-вторых, положительность квадратичного члена связывает единственность решения с выпуклостью функционала. При условии~\eqref{eq:anisotropic-ellipticity-v03} энергия строго контролирует градиент допустимой функции.

В V05 формула~\eqref{eq:poisson-variational-preview-v03} будет записана в точных функциональных пространствах и затем ограничена конечномерным пространством базисных функций. Сейчас она используется как вычислительный мост: PDE превращается в задачу о стационарной точке энергии, а после дискретизации --- в систему линейных уравнений.

\begin{caution}[энергия существует не для любого оператора]
Простая энергия~\eqref{eq:poisson-energy-v03} опирается на симметрию $A$. Для оператора переноса, несимметричной диффузии или общего нелинейного PDE аналогичный выпуклый функционал может отсутствовать. Нельзя объявлять произвольную сумму невязок ``физической энергией'' без отдельного вывода.
\end{caution}

\subsection{Законченный расчёт стационарной температуры}

\begin{problem}[стержень с равномерным тепловыделением]
Стержень длины $L=0.20\,\text{м}$ имеет постоянную теплопроводность
\[
\kappa=40\,\text{Вт}/(\text{м}\cdot\text{К}).
\]
Внутри действует равномерный источник
\[
q=8000\,\text{Вт}/\text{м}^{3},
\]
а обе границы поддерживаются при температуре $T_b=300\,\text{К}$. В одномерной модели требуется найти стационарное поле температуры и его максимум.
\end{problem}

Дано уравнение
\begin{equation}
\label{eq:rod-poisson-problem-v03}
-\kappa T''(x)=q,
\qquad
T(0)=T(L)=T_b.
\end{equation}
Делим на $-\kappa$ и дважды интегрируем:
\begin{align}
T''(x)&=-\frac{q}{\kappa},\\
T'(x)&=-\frac{q}{\kappa}x+C_1,\\
T(x)&=-\frac{q}{2\kappa}x^2+C_1x+C_2.
\end{align}
Из условия $T(0)=T_b$ получаем $C_2=T_b$. Подставляем второе граничное условие:
\[
T_b
=-\frac{q}{2\kappa}L^2+C_1L+T_b,
\]
откуда
\[
C_1=\frac{qL}{2\kappa}.
\]
Следовательно,
\begin{equation}
\label{eq:rod-temperature-solution-v03}
T(x)
=
T_b
+
\frac{q}{2\kappa}x(L-x).
\end{equation}

Максимум находится из условия $T'(x)=0$:
\[
-\frac{q}{\kappa}x+\frac{qL}{2\kappa}=0,
\qquad
x=\frac{L}{2}=0.10\,\text{м}.
\]
Подставляем численные значения:
\begin{align}
T_{\max}-T_b
&=
\frac{qL^2}{8\kappa}
=
\frac{8000\cdot(0.20)^2}{8\cdot40}
\\
&=
\frac{320}{320}
=1\,\text{К}.
\end{align}
Итак,
\begin{equation}
\label{eq:rod-temperature-maximum-v03}
T_{\max}=301\,\text{К}.
\end{equation}

Проверим баланс. Тепловой поток равен
\[
j(x)=-\kappa T'(x)
=q\left(x-\frac{L}{2}\right).
\]
На левой границе внешняя нормаль направлена влево, поэтому выходящий поток равен $800\,\text{Вт}/\text{м}^{2}$. На правой границе получаем такое же значение. Их сумма
\[
800+800=1600\,\text{Вт}/\text{м}^{2}
\]
совпадает с интегралом источника на единицу площади поперечного сечения:
\[
\int_0^L q\,\dd x=qL=1600\,\text{Вт}/\text{м}^{2}.
\]
Расчёт проверяет одновременно решение PDE, граничные условия и закон сохранения.

\subsection{Роль в генеративном проектировании и ограничения}

Для конструкции коэффициент в эллиптическом операторе зависит от геометрии и материала:
\begin{equation}
\label{eq:design-poisson-v03}
-\nabla\cdot
\bigl(A(x;\theta)\nabla u(x;\theta)\bigr)
=f(x;\theta).
\end{equation}
Параметры $\theta$ могут задавать форму области, локальную толщину, ориентацию анизотропного материала или распределение пористости. Решение $u(\cdot;\theta)$ превращает описание кандидата в поле физического отклика.

Вычислительная цепочка имеет конкретный вид
\[
\theta
\longmapsto
\bigl(D(\theta),A(\cdot;\theta),f(\cdot;\theta),g(\theta)\bigr)
\longmapsto
u(\cdot;\theta)
\longmapsto
J\bigl(u(\cdot;\theta),\theta\bigr).
\]
Для тепловой задачи функционал может измерять максимальную температуру, среднюю температуру или энергию~\eqref{eq:poisson-energy-v03}. Для поля потенциала интерес может представлять поток через выбранную часть границы.

\begin{mlconnection}[что именно проверяет физический решатель]
Генеративная модель предлагает параметры $\theta$ или поле коэффициентов $A(x;\theta)$. PDE-решатель не подтверждает конструкцию одним числом: он вычисляет распределённое поле $u$. После этого проверяются граничные условия, баланс потоков, допустимость коэффициентов и целевой функционал. Малое значение PDE-невязки в отдельных точках не заменяет эти проверки.
\end{mlconnection}

Главное ограничение связано с эллиптичностью. Если генератор выдаёт материал с нулевой или отрицательной проводимостью там, где модель предполагает~\eqref{eq:anisotropic-ellipticity-v03}, задача может потерять устойчивость или физический смысл. Простое отсечение коэффициента снизу меняет сгенерированную конструкцию и должно учитываться при обучении и оценке.

Геометрия также влияет на постановку. Узкие перемычки, острые углы и почти соприкасающиеся границы ухудшают качество сетки и могут создавать большие градиенты. Принцип максимума гарантирует свойства непрерывного решения при своих предположениях, но не исправляет ошибочную геометрию или грубую дискретизацию.

Условия Дирихле, Неймана и источник должны соответствовать одному физическому сценарию. В частности, для чистой задачи Неймана обязательна совместимость~\eqref{eq:neumann-compatibility-v03}. Генерация формы без одновременного переноса нагрузок и закреплений создаёт формально определённую, но физически другую задачу.

Раздел даёт базовый эллиптический оператор, принцип единственности и энергетическую структуру. В V04 к нему добавится время, в V05 --- конечномерная аппроксимация, а в V06 --- производная решения и функционала по параметрам конструкции.

\section{Теплопроводность, волны и разностные схемы}
\label{sec:pde-time-dependent-and-finite-differences}

\subsection{Нестационарное поле и уравнение теплопроводности}

Стационарная задача из раздела~\ref{sec:laplace-poisson-equations} описывает уже установившееся поле. Если требуется проследить нагрев, охлаждение или распространение возмущения, неизвестное поле зависит от времени:
\[
u=u(x,t),
\qquad
x\in D,
\quad
0<t\leq T.
\]
Для однородного изотропного материала уравнение теплопроводности имеет вид
\begin{equation}
\label{eq:heat-equation-v04}
 u_t-\kappa\Delta u=s,
\end{equation}
где $\kappa>0$ --- температуропроводность, а $s(x,t)$ --- объёмный источник. При неоднородном материале коэффициент входит внутрь дивергенции:
\begin{equation}
\label{eq:heterogeneous-heat-equation-v04}
 c(x)u_t
 -\nabla\cdot\bigl(k(x)\nabla u\bigr)
 =q(x,t).
\end{equation}
Здесь $c(x)>0$ задаёт объёмную теплоёмкость, $k(x)>0$ --- теплопроводность, а $q$ --- мощность источника на единицу объёма. Такое разделение важно в проектировании: изменение материала одновременно меняет накопление и перенос тепла.

Начально-краевая задача дополняется начальным полем
\begin{equation}
\label{eq:heat-initial-data-v04}
 u(x,0)=u_0(x)
\end{equation}
и граничными условиями на $\partial D$. При фиксированной температуре используется условие Дирихле, при заданном тепловом потоке --- условие Неймана, а при теплообмене со средой --- условие Робина.

Лапласиан сравнивает значение поля в точке со значениями поблизости. Поэтому локальный максимум температуры убывает, а локальный минимум возрастает. Высокочастотные пространственные колебания затухают быстрее низкочастотных. Уравнение теплопроводности сглаживает поле и ведёт его к стационарному состоянию, которое при $u_t=0$ удовлетворяет эллиптическому уравнению из V03.

Обратная эволюция принципиально несимметрична прямой. Малые высокочастотные ошибки конечного температурного поля при движении назад усиливаются. Поэтому восстановление прошлого состояния по позднему измерению является плохо обусловленной обратной задачей, даже если прямое моделирование устойчиво.

\begin{caution}[диффузия не означает произвольное размытие]
Коэффициенты, источник, начальные и граничные данные определяют конкретный процесс. Замена физического оператора на сглаживающий фильтр может визуально давать похожий результат, но не сохраняет поток, размерности и зависимость от материала. В оптимизации конструкции это приводит к градиенту другой задачи.
\end{caution}

\subsection{Волновое уравнение и энергия}

Волновое уравнение содержит вторую производную по времени:
\begin{equation}
\label{eq:wave-equation-v04}
 u_{tt}-c^2\Delta u=f,
\end{equation}
где $c>0$ --- скорость распространения волны. В отличие от теплопроводности, требуется задать два начальных поля:
\begin{equation}
\label{eq:wave-initial-data-v04}
 u(x,0)=u_0(x),
 \qquad
 u_t(x,0)=v_0(x).
\end{equation}
Первое задаёт начальное смещение, второе --- начальную скорость.

Для однородной задачи $f=0$ с закреплённой границей $u=0$ естественная энергия равна
\begin{equation}
\label{eq:wave-energy-v04}
 E(t)
 =
 \frac12\int_D
 \left(
  u_t(x,t)^2
  +c^2\norm{\nabla u(x,t)}^2
 \right)dx.
\end{equation}
Умножим~\eqref{eq:wave-equation-v04} на $u_t$ и проинтегрируем по $D$. Интегрирование пространственного члена по частям и нулевое граничное значение дают
\[
\frac{d}{dt}E(t)=0.
\]
Следовательно, энергия не исчезает, а переходит между кинетической частью $u_t^2/2$ и частью, связанной с деформацией $c^2\norm{\nabla u}^2/2$.

Теплопроводность и волна по-разному обрабатывают мелкомасштабные детали. Параболическое уравнение сглаживает их, гиперболическое переносит и отражает от границ. Поэтому численная схема для волны должна контролировать не только рост решения, но и фазовую ошибку: вычисленная волна может сохранять ограниченную амплитуду, но двигаться с неправильной скоростью.

\begin{application}[какой тип динамики нужен конструкции]
Температурное поле электронного узла обычно моделируется параболической задачей. Кратковременный удар, вибрация или акустическое возмущение требуют гиперболической модели. Подмена волновой динамики последовательностью стационарных задач теряет инерцию, скорость распространения и резонансные эффекты.
\end{application}

\subsection{Пространственно-временная сетка}

Для вычисления поля область пространства и интервал времени заменяются конечными наборами точек. В одном измерении положим
\begin{equation}
\label{eq:space-time-grid-v04}
 x_i=ih,
 \qquad
 t_n=n\tau,
\end{equation}
где $h>0$ --- шаг по пространству, $\tau>0$ --- шаг по времени. Число
\[
 u_i^n
 \approx
 u(x_i,t_n)
\]
обозначает приближённое значение поля в узле $x_i$ на временном слое $t_n$.

Производная по времени заменяется прямой разностью
\begin{equation}
\label{eq:forward-time-difference-v04}
 u_t(x_i,t_n)
 \approx
 \frac{u_i^{n+1}-u_i^n}{\tau}.
\end{equation}
Вторая пространственная производная заменяется центральной разностью
\begin{equation}
\label{eq:centered-space-difference-v04}
 u_{xx}(x_i,t_n)
 \approx
 \frac{u_{i+1}^n-2u_i^n+u_{i-1}^n}{h^2}.
\end{equation}
При достаточно гладком решении первая формула имеет ошибку порядка $\tau$, а вторая --- порядка $h^2$. Порядок показывает скорость уменьшения локальной ошибки при измельчении сетки, но сам по себе не гарантирует устойчивость всей последовательности шагов.

Граничные узлы обрабатываются в соответствии с физическими данными. Для условия Дирихле их значения задаются непосредственно. Для условия Неймана требуется разностная аппроксимация нормальной производной или дополнительный фиктивный узел. Нельзя сначала обновить все узлы одной формулой, а затем произвольно исправить границу: это меняет дискретный оператор.

В многомерной области регулярная сетка удобна только для простой геометрии. Для сложной формы конструкции используются треугольники или тетраэдры и конечные элементы, которые будут введены в V05. Однако логика остаётся той же: пространственная дискретизация превращает поле в конечномерный вектор, а эволюционное PDE --- в систему ОДУ.

\subsection{Явный и неявный временной шаг}

Подставим~\eqref{eq:forward-time-difference-v04} и~\eqref{eq:centered-space-difference-v04} в одномерное уравнение
\[
 u_t=\kappa u_{xx}+s.
\]
Получаем явную схему
\begin{equation}
\label{eq:explicit-heat-scheme-v04}
 u_i^{n+1}
 =
 (1-2r)u_i^n
 +r u_{i-1}^n
 +r u_{i+1}^n
 +\tau s_i^n,
 \qquad
 r=\frac{\kappa\tau}{h^2}.
\end{equation}
Она называется явной, потому что новый слой вычисляется непосредственно по уже известному слою.

При $s=0$ и
\begin{equation}
\label{eq:heat-cfl-v04}
 0\leq r\leq\frac12
\end{equation}
коэффициенты в~\eqref{eq:explicit-heat-scheme-v04} неотрицательны и суммируются в единицу. Новое значение является выпуклой комбинацией трёх старых значений. Следовательно, схема не создаёт новый максимум или минимум внутри области. Это дискретный аналог принципа максимума.

То же ограничение видно из анализа отдельной гармоники. Для ошибки вида
\[
 e_i^n=G^n e^{\mathrm{i}\xi x_i}
\]
коэффициент усиления равен
\begin{equation}
\label{eq:heat-amplification-factor-v04}
 G(\xi)
 =
 1-4r\sin^2\!\left(\frac{\xi h}{2}\right).
\end{equation}
Условие $\abs{G(\xi)}\leq1$ для всех пространственных частот снова приводит к $r\leq1/2$. Если шаг времени слишком велик, некоторые сеточные колебания меняют знак и растут, хотя непрерывное уравнение должно их подавлять.

Неявная схема использует пространственную производную на новом слое:
\begin{equation}
\label{eq:implicit-heat-scheme-v04}
 u_i^{n+1}
 -r\left(
 u_{i+1}^{n+1}-2u_i^{n+1}+u_{i-1}^{n+1}
 \right)
 =
 u_i^n+\tau s_i^{n+1}.
\end{equation}
Теперь на каждом шаге решается линейная система. Такой шаг дороже явного, но для линейной теплопроводности не имеет ограничения~\eqref{eq:heat-cfl-v04} того же типа. Большой $\tau$ всё равно снижает точность и может скрыть быстрые переходные процессы.

После конечных элементов временная задача принимает матричный вид
\begin{equation}
\label{eq:semidiscrete-parabolic-system-v04}
 M\dot U(t)+K(t)U(t)=F(t),
\end{equation}
где $M$ --- матрица масс, $K$ --- матрица жёсткости, а $U(t)$ --- коэффициенты поля. Неявный Эйлер даёт систему
\begin{equation}
\label{eq:implicit-euler-matrix-v04}
 \bigl(M+\tau K_{n+1}\bigr)U^{n+1}
 =
 MU^n+\tau F^{n+1}.
\end{equation}
Так непрерывная эволюция сводится к последовательности стационарных линейных задач.

\subsection{Законченный расчёт одного шага теплопроводности}

\begin{problem}[законченный расчёт явного шага]
Одномерный стержень разбит узлами с шагом
\[
 h=0.01\,\text{м}.
\]
Температуропроводность равна
\[
 \kappa=10^{-4}\,\text{м}^2/\text{с},
\]
шаг времени --- $\tau=0.4\,\text{с}$, внутреннего источника нет. На текущем слое температуры равны
\begin{equation}
\label{eq:heat-step-data-v04}
 \bigl(u_0^n,u_1^n,u_2^n,u_3^n,u_4^n\bigr)
 =
 (300,320,360,320,300)\,\text{К},
\end{equation}
а крайние узлы поддерживаются при $300\,\text{К}$. Найдём следующий временной слой по явной схеме.
\end{problem}

Сначала вычисляем безразмерный параметр:
\begin{align}
 r
 &=
 \frac{\kappa\tau}{h^2}
 =
 \frac{10^{-4}\cdot0.4}{(0.01)^2}
 =0.4.
\end{align}
Условие устойчивости выполнено, поскольку $0.4<1/2$. Для первого внутреннего узла
\begin{align}
 u_1^{n+1}
 &=(1-2r)u_1^n+r u_0^n+r u_2^n
 \\
 &=0.2\cdot320+0.4\cdot300+0.4\cdot360
 \\
 &=64+120+144
 =328\,\text{К}.
\end{align}
Для центрального узла
\begin{align}
 u_2^{n+1}
 &=0.2\cdot360+0.4\cdot320+0.4\cdot320
 \\
 &=72+128+128
 =328\,\text{К}.
\end{align}
По симметрии $u_3^{n+1}=328\,\text{К}$. Граничные значения не обновляются, поэтому
\begin{equation}
\label{eq:heat-step-result-v04}
 \bigl(u_0^{n+1},u_1^{n+1},u_2^{n+1},u_3^{n+1},u_4^{n+1}\bigr)
 =
 (300,328,328,328,300)\,\text{К}.
\end{equation}
Пик $360\,\text{К}$ уменьшился, а соседние узлы нагрелись. Это соответствует переносу тепла от более горячего центра к более холодным частям.

Если оставить ту же сетку, но взять $\tau=0.6\,\text{с}$, то $r=0.6$ и центральный вес $1-2r=-0.2$ становится отрицательным. Один шаг ещё может выглядеть правдоподобно для специально выбранных данных, но схема больше не сохраняет дискретный принцип максимума и не является устойчивой для произвольных сеточных возмущений.

\subsection{Временной решатель в ML и генеративном проектировании}

Пусть $\theta$ задаёт форму конструкции, распределение материала или параметры источника. После пространственной дискретизации вычислительная модель имеет вид
\begin{equation}
\label{eq:design-time-system-v04}
 M(\theta)\dot U(t;\theta)
 +K(\theta,t)U(t;\theta)
 =F(\theta,t).
\end{equation}
Решатель строит последовательность состояний
\[
 U^0(\theta),U^1(\theta),\ldots,U^N(\theta).
\]
Целевой функционал может учитывать конечное состояние, всю траекторию или худший момент времени:
\begin{equation}
\label{eq:transient-objective-v04}
 J(\theta)
 =
 \Phi\bigl(U^N(\theta),\theta\bigr)
 +\tau\sum_{n=0}^{N}
 \ell\bigl(U^n(\theta),\theta,t_n\bigr).
\end{equation}
Так можно штрафовать перегрев, длительное превышение допустимой температуры, вибрации или накопленную энергию.

\begin{mlconnection}[дифференцируемый временной шаг]
Явная формула~\eqref{eq:explicit-heat-scheme-v04} является композицией арифметических операций и потому легко дифференцируется автоматическим дифференцированием. Однако уменьшение $h$ требует уменьшать $\tau$ как $h^2$, поэтому число шагов быстро растёт. Неявный шаг~\eqref{eq:implicit-euler-matrix-v04} требует дифференцировать решение линейной системы, но допускает более крупный шаг времени. Выбор схемы определяет одновременно стоимость прямого расчёта, память обратного прохода и качество градиента.
\end{mlconnection}

Связь со Сбером и AIRI здесь состоит не в добавлении нейросети внутрь PDE, а в построении проверяемого вычислительного конвейера. Генеративная модель предлагает параметр $\theta$, физический решатель вычисляет переходный отклик, а критерий принимает или отклоняет конструкцию. Суррогатная ML-модель может ускорять отдельный этап, но должна проверяться относительно исходной дискретной задачи на геометриях, материалах и временных режимах, встречающихся при оптимизации.

Нужно различать ошибку модели, пространственную ошибку, временную ошибку и ошибку оптимизации. Уменьшение шага времени не исправляет неверные граничные условия. Точная линейная система не исправляет грубую сетку. Малый тренировочный loss суррогата не гарантирует сохранение энергии волны или дискретного принципа максимума теплопроводности.

При изменении топологии сетка и размерность вектора $U$ могут меняться скачком. Тогда простое дифференцирование через последовательность матриц перестаёт описывать производную одной фиксированной вычислительной схемы. Этот вопрос будет отделён от V04: V05 введёт конечные элементы и сборку матриц, а V06 разберёт производную решения при фиксированной дискретной структуре и явно укажет границы применимости.

\section{Слабая постановка и метод конечных элементов}
\label{sec:pde-weak-form-and-fem}

\subsection{Зачем ослаблять понятие решения}

Классическое решение эллиптического уравнения второго порядка должно иметь вторые производные. Для реальной конструкции это требование часто слишком сильное: коэффициент материала может скачком меняться на границе двух фаз, область может иметь углы, а нагрузка может быть задана лишь интегрируемой функцией. Поле при этом остаётся физически осмысленным, хотя его вторые производные не существуют в каждой точке.

Слабая постановка переносит одну производную с неизвестного поля на пробную функцию. Поэтому для уравнения второго порядка достаточно контролировать первые производные решения. Нужный класс задаётся пространством
\begin{equation}
\label{eq:sobolev-h1-v05}
 H^1(D)
 =
 \left\{
 u\in L^2(D)
 \,\middle|\,
 \partial_i u\in L^2(D),\ i=1,\ldots,d
 \right\},
\end{equation}
где производные понимаются в слабом смысле.

\begin{definition}
Функция $g_i\in L^2(D)$ называется слабой производной $\partial_i u$, если
\begin{equation}
\label{eq:weak-derivative-v05}
 \int_D u\,\partial_i\varphi\,\dd x
 =
 -\int_D g_i\varphi\,\dd x
\end{equation}
для любой гладкой функции $\varphi$ с компактным носителем внутри $D$.
\end{definition}

Если $u$ гладкая, равенство~\eqref{eq:weak-derivative-v05} совпадает с обычным интегрированием по частям. Определение важно тем, что не требует вычислять производную в каждой точке. Например, непрерывная кусочно-линейная функция имеет постоянную производную на каждом отрезке и принадлежит $H^1$, хотя её классическая производная скачком меняется в узлах.

Однородное условие Дирихле кодируется пространством
\begin{equation}
\label{eq:h10-space-v05}
 H_0^1(D)
 =
 \overline{C_c^\infty(D)}^{\,H^1}.
\end{equation}
Практически это означает нулевое граничное значение в смысле следа. Полная теория следов здесь не нужна: при конечных элементах условие реализуется фиксацией узловых значений на границе.

\begin{connection}[связь с предыдущими главами]
Пространство $H^1(D)$ является бесконечномерным аналогом евклидова пространства. Вместо суммы квадратов координат используется интеграл от квадрата функции и её градиента. Слабая постановка затем превращает PDE в линейное уравнение в этом пространстве, а метод конечных элементов выбирает конечномерное подпространство и возвращает задачу к матричной линейной алгебре главы~\ref{chap:linear-algebra}.
\end{connection}

\subsection{Переход от сильной формы к слабой}

Рассмотрим задачу Пуассона с однородным условием Дирихле:
\begin{equation}
\label{eq:poisson-strong-v05}
 -\nabla\cdot\bigl(A(x)\nabla u(x)\bigr)=f(x)
 \quad\text{в }D,
 \qquad
 u=0
 \quad\text{на }\partial D.
\end{equation}
Умножим уравнение на пробную функцию $v\in H_0^1(D)$ и проинтегрируем. Формула Грина даёт
\begin{align}
 \int_D
 -\nabla\cdot(A\nabla u)\,v\,\dd x
 &=
 \int_D
 (A\nabla u)^{\trans}\nabla v\,\dd x
 -
 \int_{\partial D}
 v\,n^{\trans}A\nabla u\,\dd s.
\end{align}
Граничный интеграл исчезает, поскольку $v=0$ на $\partial D$. Получаем слабую постановку:
\begin{equation}
\label{eq:poisson-weak-v05}
 \text{найти }u\in H_0^1(D):
 \qquad
 a(u,v)=\ell(v)
 \quad
 \text{для всех }v\in H_0^1(D),
\end{equation}
где
\begin{equation}
\label{eq:bilinear-and-load-v05}
 a(u,v)
 =
 \int_D (A\nabla u)^{\trans}\nabla v\,\dd x,
 \qquad
 \ell(v)
 =
 \int_D fv\,\dd x.
\end{equation}

Слабая форма не является приближением PDE. Она задаёт другое, менее требовательное понятие решения. Если слабое решение оказывается достаточно гладким, то оно удовлетворяет исходному уравнению почти всюду и совпадает с классическим решением.

Условия Дирихле входят в пространство допустимых функций и называются существенными. Условия Неймана появляются после интегрирования по частям в граничном интеграле и называются естественными. Для смешанной задачи с $u=0$ на $\Gamma_D$ и $n^{\trans}A\nabla u=g_N$ на $\Gamma_N$ правая часть становится
\begin{equation}
\label{eq:weak-neumann-load-v05}
 \ell(v)
 =
 \int_D fv\,\dd x
 +
 \int_{\Gamma_N}g_Nv\,\dd s.
\end{equation}

\begin{caution}[знак граничного потока]
Знак в~\eqref{eq:weak-neumann-load-v05} зависит от того, как в сильной форме определён поток и какая нормаль выбрана внешней. Ошибка знака не является малой численной погрешностью: она меняет направление приложенной нагрузки или теплового потока.
\end{caution}

\subsection{Корректность и энергетический смысл}

Абстрактная слабая задача имеет вид
\begin{equation}
\label{eq:abstract-weak-problem-v05}
 \text{найти }u\in V:
 \qquad
 a(u,v)=\ell(v)
 \quad
 \text{для всех }v\in V,
\end{equation}
где $V$ --- гильбертово пространство. Для надёжности постановки нужны два свойства билинейной формы. Непрерывность означает существование $M>0$ такого, что
\begin{equation}
\label{eq:bilinear-continuity-v05}
 \abs{a(u,v)}
 \leq
 M\norm{u}_V\norm{v}_V.
\end{equation}
Коэрцитивность означает существование $\alpha>0$ такого, что
\begin{equation}
\label{eq:bilinear-coercivity-v05}
 a(v,v)
 \geq
 \alpha\norm{v}_V^2.
\end{equation}

\begin{theorem}[Лакса--Мильграма]
Если $a$ непрерывна и коэрцитивна на $V$, а $\ell$ является непрерывным линейным функционалом, то задача~\eqref{eq:abstract-weak-problem-v05} имеет единственное решение. Кроме того,
\begin{equation}
\label{eq:lax-milgram-stability-v05}
 \norm{u}_V
 \leq
 \frac{1}{\alpha}\norm{\ell}_{V^*}.
\end{equation}
\end{theorem}

Для симметричной формы слабое решение одновременно минимизирует энергию
\begin{equation}
\label{eq:weak-energy-v05}
 \mathcal E(v)
 =
 \frac12 a(v,v)-\ell(v).
\end{equation}
Действительно, производная по направлению $w$ равна
\[
 \fdv{\mathcal E(v)}{v}[w]
 =a(v,w)-\ell(w),
\]
поэтому условие стационарности для всех $w\in V$ совпадает со слабой задачей.

Для генеративного проектирования это различие важно. PDE задаёт допустимый физический отклик, а функционал~\eqref{eq:weak-energy-v05} может быть внутренней энергией состояния. Его нельзя автоматически отождествлять с целевой функцией проектирования: задача оптимизации может дополнительно штрафовать массу, максимум температуры, ограничение напряжений или технологические свойства геометрии.

\subsection{Метод Галёркина и матричная система}

Выберем конечномерное подпространство
\[
 V_h=\spanop\{\varphi_1,\ldots,\varphi_N\}
 \subset V.
\]
Метод Галёркина ищет $u_h\in V_h$ такое, что
\begin{equation}
\label{eq:galerkin-problem-v05}
 a(u_h,v_h)=\ell(v_h)
 \qquad
 \text{для всех }v_h\in V_h.
\end{equation}
Разложим решение по базису:
\begin{equation}
\label{eq:galerkin-expansion-v05}
 u_h(x)=\sum_{j=1}^{N}U_j\varphi_j(x).
\end{equation}
Достаточно подставить в~\eqref{eq:galerkin-problem-v05} базисные функции $v_h=\varphi_i$. Получаем
\begin{equation}
\label{eq:fem-linear-system-v05}
 KU=F,
 \qquad
 K_{ij}=a(\varphi_j,\varphi_i),
 \qquad
 F_i=\ell(\varphi_i).
\end{equation}
Матрица $K$ называется матрицей жёсткости. Для нестационарных задач также возникает матрица масс
\begin{equation}
\label{eq:fem-mass-matrix-v05}
 M_{ij}=\int_D\varphi_j\varphi_i\,\dd x,
\end{equation}
что согласуется с системой $M\dot U+KU=F$ из V04.

Из непрерывности и коэрцитивности следует оценка квазинаилучшего приближения
\begin{equation}
\label{eq:cea-estimate-v05}
 \norm{u-u_h}_V
 \leq
 \frac{M}{\alpha}
 \inf_{v_h\in V_h}\norm{u-v_h}_V.
\end{equation}
Она отделяет две причины ошибки. Свойства PDE входят через отношение $M/\alpha$, а качество пространства $V_h$ --- через способность его функций приближать точное решение. Увеличение числа узлов не помогает, если задача почти вырождена и $\alpha$ близка к нулю.

\subsection{Конечные элементы и законченная сборка системы}

Метод конечных элементов строит $V_h$ из локальных полиномиальных функций на сетке. В одном измерении область разбивается на отрезки, в двух --- обычно на треугольники или четырёхугольники, в трёх --- на тетраэдры или другие ячейки. Непрерывные кусочно-линейные базисные функции $\varphi_i$ равны единице в своём узле и нулю во всех остальных узлах. Их носители локальны, поэтому большинство элементов $K_{ij}$ равно нулю и матрица получается разреженной.

Сборка выполняется по элементам. На каждом элементе $T$ вычисляются локальные матрица и вектор
\begin{equation}
\label{eq:local-fem-contributions-v05}
 K^{(T)}_{ab}
 =
 \int_T
 (A\nabla\varphi_b)^{\trans}\nabla\varphi_a\,\dd x,
 \qquad
 F^{(T)}_a
 =
 \int_T f\varphi_a\,\dd x.
\end{equation}
Затем локальные номера узлов переводятся в глобальные, а вклады складываются в общие $K$ и $F$. Один узел принадлежит нескольким элементам, поэтому соответствующая строка глобальной системы получает сумму всех соседних вкладов.

\begin{problem}[законченная сборка малой FEM-системы]
Решим задачу
\begin{equation}
\label{eq:fem-worked-problem-v05}
 -u''(x)=1,
 \qquad
 x\in(0,1),
 \qquad
 u(0)=u(1)=0
\end{equation}
кусочно-линейными конечными элементами на равномерной сетке
\[
 x_0=0,
 \qquad
 x_1=\frac13,
 \qquad
 x_2=\frac23,
 \qquad
 x_3=1.
\]
\end{problem}

Длина каждого элемента равна $h=1/3$. На элементе $T_e=[x_{e-1},x_e]$ производные двух локальных базисных функций равны $-1/h$ и $1/h$. Поэтому
\begin{equation}
\label{eq:local-stiffness-worked-v05}
 K^{(e)}
 =
 \frac1h
 \begin{pmatrix}
 1&-1\\
 -1&1
 \end{pmatrix}
 =
 3
 \begin{pmatrix}
 1&-1\\
 -1&1
 \end{pmatrix}.
\end{equation}
При $f=1$ локальный вектор нагрузки равен
\begin{equation}
\label{eq:local-load-worked-v05}
 F^{(e)}
 =
 \frac h2
 \begin{pmatrix}
 1\\
 1
 \end{pmatrix}
 =
 \frac16
 \begin{pmatrix}
 1\\
 1
 \end{pmatrix}.
\end{equation}

Складывая вклады трёх элементов, получаем глобальную систему до учёта условий Дирихле:
\begin{equation}
\label{eq:global-fem-system-v05}
 \begin{pmatrix}
 3&-3&0&0\\
 -3&6&-3&0\\
 0&-3&6&-3\\
 0&0&-3&3
 \end{pmatrix}
 \begin{pmatrix}
 U_0\\U_1\\U_2\\U_3
 \end{pmatrix}
 =
 \begin{pmatrix}
 1/6\\1/3\\1/3\\1/6
 \end{pmatrix}.
\end{equation}
Граничные значения известны: $U_0=U_3=0$. Для двух внутренних неизвестных остаётся
\begin{equation}
\label{eq:reduced-fem-system-v05}
 \begin{pmatrix}
 6&-3\\
 -3&6
 \end{pmatrix}
 \begin{pmatrix}
 U_1\\U_2
 \end{pmatrix}
 =
 \begin{pmatrix}
 1/3\\1/3
 \end{pmatrix}.
\end{equation}
Из симметрии $U_1=U_2=:q$, поэтому $3q=1/3$ и
\begin{equation}
\label{eq:fem-worked-solution-v05}
 U_1=U_2=\frac19.
\end{equation}
Искомое поле равно
\[
 u_h(x)=\frac19\varphi_1(x)+\frac19\varphi_2(x).
\]
Точное решение задачи~\eqref{eq:fem-worked-problem-v05} имеет вид $u(x)=x(1-x)/2$ и в узлах $x_1,x_2$ также равно $1/9$. Между узлами конечный элемент заменяет параболу линейными отрезками, поэтому совпадение узловых значений не означает нулевую ошибку поля.

\subsection{FEM в генеративном проектировании и границы применимости}

После построения сетки вычислительная модель принимает вид
\begin{equation}
\label{eq:fem-design-system-v05}
 K(\theta)U(\theta)=F(\theta),
\end{equation}
где $\theta$ может задавать коэффициенты материала, положение границы, толщины или параметры геометрии. Это ключевой интерфейс между непрерывной физикой и ML: модель работает с конечномерными тензорами, но их значения получены из интегралов по реальной области и сохраняют смысл исходного PDE.

\begin{mlconnection}[что именно можно обучать]
Суррогатная модель может приближать отображение $\theta\mapsto U(\theta)$, предсказывать локальную ошибку сетки или предлагать начальное приближение для линейного решателя. Генеративная модель может создавать поле материала или геометрию. Во всех случаях физическая проверка должна возвращаться к системе~\eqref{eq:fem-design-system-v05} и к функционалу, вычисленному по восстановленному полю, а не только к похожести предсказания на тренировочные примеры.
\end{mlconnection}

Качество FEM определяется не одним числом элементов. Вытянутые и почти вырожденные ячейки ухудшают обусловленность. Разрывные коэффициенты требуют сетки или метода, способных корректно представить интерфейс. Сингулярности в углах замедляют сходимость равномерного сгущения. Численное интегрирование низкого порядка может неверно собрать матрицу даже при хорошей сетке.

При изменении топологии конструкции меняются сетка, связи между узлами и иногда размерность $U$. Тогда матрица $K(\theta)$ не является гладкой функцией параметра в обычном смысле на всём пространстве проектов. Кроме того, малая алгебраическая невязка $\norm{KU-F}$ означает лишь точное решение выбранной дискретной системы; она не контролирует ошибку физической модели и ошибку дискретизации.

V05 завершает переход
\[
 \text{PDE}
 \longrightarrow
 \text{слабая форма}
 \longrightarrow
 \text{конечномерное пространство}
 \longrightarrow
 KU=F.
\]
В V06 эта система будет рассматриваться как неявно заданное отображение параметров конструкции в поле отклика. Мы выведем производную $U(\theta)$ и сопряжённую формулу для градиента целевого функционала при фиксированной дискретной структуре.

% V06
% \section{Дифференцируемый PDE-решатель и оптимизация конструкции}

\section{Дифференцируемый PDE-решатель и градиент по конструкции}
\label{sec:differentiable-pde-solver}

\subsection{Параметрическая задача и отображение решения}

До сих пор коэффициенты, область и нагрузки считались заданными. В проектировании они зависят от параметров
\[
 \theta=(\theta_1,\ldots,\theta_p)\in\Theta\subset\R^p.
\]
Параметр может задавать толщину, локальную плотность материала, теплопроводность, форму границы или значение нагрузки. После дискретизации слабой задачи возникает система
\begin{equation}
\label{eq:parametric-pde-system-v06}
 K(\theta)U(\theta)=F(\theta).
\end{equation}
Здесь $U(\theta)\in\R^N$ содержит коэффициенты поля, $K(\theta)$ является дискретным оператором, а $F(\theta)$ --- вектором данных. Решатель реализует отображение
\begin{equation}
\label{eq:pde-solution-map-v06}
 \mathcal S_h:\theta\longmapsto U(\theta).
\end{equation}

Для заданного поля вычисляется целевая функция
\begin{equation}
\label{eq:reduced-objective-v06}
 j(\theta)
 =
 \Phi\bigl(U(\theta),\theta\bigr).
\end{equation}
Она называется приведённой, поскольку зависимость от состояния уже подчинена уравнению~\eqref{eq:parametric-pde-system-v06}. Именно производная $\nabla_\theta j$ нужна алгоритму оптимизации или обучаемой модели, предлагающей параметры конструкции.

Существование решения системы ещё не гарантирует существование устойчивой производной. Для локальной дифференцируемости достаточно, чтобы $K$ и $F$ дифференцировались по $\theta$, а матрица $K(\theta)$ оставалась обратимой в рассматриваемой окрестности. Если минимальное собственное или сингулярное число приближается к нулю, малое изменение параметра может вызвать большое изменение поля. Это дискретный аналог потери устойчивости краевой задачи.

В слабой форме та же конструкция записывается как
\begin{equation}
\label{eq:parametric-weak-problem-v06}
 a_\theta(u_\theta,v)
 =
 \ell_\theta(v)
 \qquad
 \text{для всех }v\in V.
\end{equation}
Параметр входит в билинейную форму, правую часть, а иногда и в пространство допустимых функций. В этой главе подробно рассматривается наиболее прозрачный случай фиксированной области и фиксированного пространства $V$. Производные по движущейся границе требуют отдельного аппарата и не сводятся к простому дифференцированию коэффициентов матрицы.

\subsection{Неявное дифференцирование дискретной системы}

Равенство~\eqref{eq:parametric-pde-system-v06} определяет $U$ не явной формулой, а условием. Дифференцируя его по одному параметру $\theta_i$, получаем
\begin{equation}
\label{eq:direct-sensitivity-derivation-v06}
 \pdv{K}{\theta_i}U
 +
 K\pdv{U}{\theta_i}
 =
 \pdv{F}{\theta_i}.
\end{equation}
Следовательно, чувствительность состояния $S_i=\partial U/\partial\theta_i$ находится из системы
\begin{equation}
\label{eq:direct-sensitivity-system-v06}
 KS_i
 =
 \pdv{F}{\theta_i}
 -
 \pdv{K}{\theta_i}U.
\end{equation}
Матрица здесь та же, что и в прямой задаче. Если её факторизация уже построена, её можно повторно использовать для всех правых частей.

Полная производная целевой функции равна
\begin{equation}
\label{eq:total-objective-derivative-v06}
 \pdv{j}{\theta_i}
 =
 \pdv{\Phi}{\theta_i}
 +
 \left(\nabla_U\Phi\right)^{\trans}S_i.
\end{equation}
Первое слагаемое учитывает явную зависимость критерия от параметра, второе --- изменение критерия через физическое поле.

Прямой метод естественен, когда параметров мало, а выходных величин много. Для каждого $\theta_i$ требуется решить одну систему~\eqref{eq:direct-sensitivity-system-v06}. При миллионах переменных поля и нескольких параметрах это приемлемо. В генеративном проектировании ситуация часто обратная: параметров материала или геометрии очень много, но целевая функция одна. Тогда вычислять все $S_i$ невыгодно.

На уровне слабой формы чувствительность $s_i=\partial u_\theta/\partial\theta_i$ удовлетворяет
\begin{equation}
\label{eq:weak-direct-sensitivity-v06}
 a_\theta(s_i,v)
 =
 \pdv{}{\theta_i}\ell_\theta(v)
 -
 \pdv{}{\theta_i}a_\theta(u_\theta,v)
 \qquad
 \text{для всех }v\in V.
\end{equation}
Это равенство показывает, что чувствительность сама является решением PDE с тем же главным оператором и новой правой частью.

\subsection{Сопряжённая задача}

Сопряжённый метод устраняет необходимость явно строить все чувствительности. Введём вектор $\lambda\in\R^N$ как решение
\begin{equation}
\label{eq:discrete-adjoint-v06}
 K(\theta)^{\trans}\lambda
 =
 \nabla_U\Phi\bigl(U(\theta),\theta\bigr).
\end{equation}
Умножая~\eqref{eq:direct-sensitivity-system-v06} слева на $\lambda^{\trans}$ и используя~\eqref{eq:discrete-adjoint-v06}, получаем
\begin{equation}
\label{eq:adjoint-gradient-v06}
 \pdv{j}{\theta_i}
 =
 \pdv{\Phi}{\theta_i}
 +
 \lambda^{\trans}
 \left(
  \pdv{F}{\theta_i}
  -
  \pdv{K}{\theta_i}U
 \right).
\end{equation}
После одной прямой и одной сопряжённой задачи каждая компонента градиента вычисляется через матрично-векторные произведения и локальные производные оператора.

Если $K$ симметрична, то сопряжённая система использует ту же матрицу. Это типично для многих эллиптических задач без конвективных членов. Для несимметричного оператора транспонирование существенно: подмена $K^{\trans}$ на $K$ даёт неверный градиент.

В слабой постановке сопряжённое состояние $p\in V$ определяется равенством
\begin{equation}
\label{eq:weak-adjoint-v06}
 a_\theta(v,p)
 =
 \fdv{\Phi}{u}[v]
 \qquad
 \text{для всех }v\in V.
\end{equation}
Тогда производная приведённого функционала имеет вид
\begin{equation}
\label{eq:weak-adjoint-gradient-v06}
 \pdv{j}{\theta_i}
 =
 \pdv{\Phi}{\theta_i}
 +
 \pdv{}{\theta_i}\ell_\theta(p)
 -
 \pdv{}{\theta_i}a_\theta(u_\theta,p).
\end{equation}
Формула не требует вычислять $s_i$. Она также показывает, откуда в программной реализации возникают локальные вклады градиента: параметр влияет на интегралы, из которых собираются $K$ и $F$.

\begin{application}[градиент по полю материала]
Пусть коэффициент аппроксимирован по элементам сетки и
\begin{equation}
\label{eq:elementwise-stiffness-v06}
 K(\theta)
 =
 \sum_{e=1}^{E}a_e(\theta)K^{(e)}.
\end{equation}
Тогда
\begin{equation}
\label{eq:elementwise-stiffness-derivative-v06}
 \pdv{K}{\theta_i}
 =
 \sum_{e=1}^{E}
 \pdv{a_e}{\theta_i}K^{(e)}.
\end{equation}
Производная собирается тем же поэлементным способом, что и сама матрица. Для критерия податливости $\Phi(U)=F^{\trans}U$, фиксированной нагрузки и симметричной $K$ сопряжённое состояние совпадает с прямым: $\lambda=U$. Поэтому
\begin{equation}
\label{eq:compliance-gradient-v06}
 \pdv{j}{\theta_i}
 =
 -U^{\trans}
 \pdv{K}{\theta_i}U.
\end{equation}
Знак имеет физический смысл: увеличение жёсткости в элементе, где велика локальная энергия, уменьшает податливость конструкции.
\end{application}

\subsection{Законченный расчёт градиента}

\begin{problem}[прямая и сопряжённая производные одной конструкции]
Рассмотрим параметрическую систему
\begin{equation}
\label{eq:worked-parametric-system-v06}
 \begin{pmatrix}
 2+\theta&-1\\
 -1&2
 \end{pmatrix}
 \begin{pmatrix}
 U_1\\U_2
 \end{pmatrix}
 =
 \begin{pmatrix}
 1\\0
 \end{pmatrix}
\end{equation}
и критерий
\begin{equation}
\label{eq:worked-objective-v06}
 j(\theta)
 =
 \frac12\left(U_2(\theta)-\frac12\right)^2
 +
 \frac{1}{20}\theta^2.
\end{equation}
Вычислим $j'(1)$ прямым и сопряжённым способами.
\end{problem}

При $\theta=1$ имеем
\[
 K=
 \begin{pmatrix}
 3&-1\\-1&2
 \end{pmatrix},
 \qquad
 K^{-1}
 =
 \frac15
 \begin{pmatrix}
 2&1\\1&3
 \end{pmatrix}.
\]
Поэтому прямое состояние равно
\begin{equation}
\label{eq:worked-state-v06}
 U
 =
 K^{-1}F
 =
 \begin{pmatrix}
 2/5\\1/5
 \end{pmatrix}.
\end{equation}
Производные матрицы и правой части имеют вид
\[
 K_\theta
 =
 \begin{pmatrix}
 1&0\\0&0
 \end{pmatrix},
 \qquad
 F_\theta=0.
\]
Из~\eqref{eq:direct-sensitivity-system-v06}
\begin{equation}
\label{eq:worked-direct-sensitivity-v06}
 KS
 =
 -K_\theta U
 =
 \begin{pmatrix}
 -2/5\\0
 \end{pmatrix},
 \qquad
 S
 =
 \begin{pmatrix}
 -4/25\\-2/25
 \end{pmatrix}.
\end{equation}
Явная производная штрафа равна $\Phi_\theta=\theta/10=1/10$, а
\[
 \nabla_U\Phi
 =
 \begin{pmatrix}
 0\\U_2-1/2
 \end{pmatrix}
 =
 \begin{pmatrix}
 0\\-3/10
 \end{pmatrix}.
\]
Следовательно,
\begin{equation}
\label{eq:worked-direct-gradient-v06}
 j'(1)
 =
 \frac1{10}
 +
 \begin{pmatrix}0&-3/10\end{pmatrix}
 \begin{pmatrix}-4/25\\-2/25\end{pmatrix}
 =
 \frac{31}{250}
 =0.124.
\end{equation}

Теперь решим сопряжённую систему
\begin{equation}
\label{eq:worked-adjoint-v06}
 K^{\trans}\lambda
 =
 \nabla_U\Phi,
 \qquad
 \lambda
 =
 \begin{pmatrix}
 -3/50\\-9/50
 \end{pmatrix}.
\end{equation}
По формуле~\eqref{eq:adjoint-gradient-v06}
\begin{equation}
\label{eq:worked-adjoint-gradient-v06}
 j'(1)
 =
 \frac1{10}
 +
 \lambda^{\trans}(-K_\theta U)
 =
 \frac1{10}
 +
 \begin{pmatrix}-3/50&-9/50\end{pmatrix}
 \begin{pmatrix}-2/5\\0\end{pmatrix}
 =
 \frac{31}{250}.
\end{equation}
Оба способа дают один результат. Проверка возможна и без линейной алгебры: из~\eqref{eq:worked-parametric-system-v06}
\[
 U_2(\theta)=\frac{1}{3+2\theta},
 \qquad
 U_2'(1)=-\frac{2}{25},
\]
после чего непосредственное дифференцирование~\eqref{eq:worked-objective-v06} снова даёт $31/250$.

\subsection{Дифференцируемый вычислительный граф}

Практический конвейер начинается не с команды автоматического дифференцирования, а с явного описания зависимостей:
\[
 \theta
 \longmapsto
 \bigl(D,a,f,g\bigr)
 \longmapsto
 \bigl(K(\theta),F(\theta)\bigr)
 \longmapsto
 U(\theta)
 \longmapsto
 \Phi(U,\theta).
\]
Каждая стрелка должна иметь математически определённую производную или осознанную аппроксимацию. Если генеративная модель выдаёт поле материала, оно сначала переводится в физические коэффициенты. Если она выдаёт геометрию, нужны построение сетки, перенос граничных меток и сборка оператора. Решение линейной системы является лишь одной частью графа.

Существуют два основных способа реализации. Первый сохраняет операции итерационного или прямого решателя и применяет обратное автоматическое дифференцирование к ним. Тогда градиент относится к реально выполненному конечному алгоритму. Второй использует неявную формулу~\eqref{eq:adjoint-gradient-v06} и решает отдельную сопряжённую систему. Тогда не требуется хранить все итерации прямого решателя, но необходимо согласованно вычислять производные $K$, $F$ и целевого функционала.

Для больших разреженных систем важнее не строить полные якобианы, а уметь вычислять произведения
\[
 \pdv{K}{\theta_i}U,
 \qquad
 \lambda^{\trans}\pdv{K}{\theta_i}U.
\]
Поэлементная структура FEM позволяет получать эти величины локально. Поэтому размер градиента определяется числом проектных параметров, но число глобальных PDE-решений для одного скалярного критерия остаётся равным двум: прямому и сопряжённому.

\begin{mlconnection}[PDE-решатель как слой модели]
В машинном обучении входом слоя может быть параметрическое поле материала, выходом --- температура, смещение или несколько интегральных характеристик. Обучаемые параметры обновляются по градиенту, проходящему через физический решатель. Такой слой не заменяет PDE нейросетью: он сохраняет дискретную физическую задачу и предоставляет её производную остальной модели.
\end{mlconnection}

Для генеративного проектирования особенно важен разрыв между предложением и проверкой. Генератор создаёт кандидата, PDE-решатель вычисляет отклик, а целевая функция измеряет физическое качество. Градиент показывает локальное изменение критерия при малом изменении кандидата. Он пригоден для уточнения конструкции, для обучения генератора по физическому сигналу и для анализа чувствительности, но сам по себе не гарантирует глобально лучшую форму.

\subsection{Границы применимости и итог главы}

Корректный градиент определяется относительно конкретной дискретной задачи. Изменение сетки, квадратурной формулы, допуска линейного решателя или способа задания границы меняет вычисляемое отображение. Поэтому проверять нужно не только невязку прямой системы, но и производную. Практическая проверка состоит в сравнении направленной производной
\begin{equation}
\label{eq:directional-gradient-check-v06}
 \nabla j(\theta)^{\trans}d
\end{equation}
с симметричной разностью
\begin{equation}
\label{eq:finite-difference-gradient-check-v06}
 \frac{j(\theta+\varepsilon d)-j(\theta-\varepsilon d)}{2\varepsilon}.
\end{equation}
При слишком большом $\varepsilon$ видна нелинейность, при слишком малом --- ошибки округления и неточность решения системы. Совпадение на диапазоне шагов является диагностикой реализации, но не доказательством физической адекватности модели.

Градиент может не существовать в точках, где меняется топология сетки, активируется контакт, переключается кусочно-заданный материал или возникает разрыв из-за пороговой обработки. Сглаживание таких операций делает оптимизацию удобнее, но меняет исходную задачу. Это изменение должно быть записано как часть модели, а не скрыто внутри программного кода.

Плохая обусловленность усиливает как ошибку состояния, так и ошибку чувствительности. Для итерационного решателя условие остановки прямой задачи должно быть достаточно строгим относительно требуемой точности градиента. Слишком грубое решение может дать правдоподобное поле, но неверное направление оптимизации.

В результате глава проводит полный путь от поля к производной:
\[
\begin{aligned}
 \text{PDE}
 &\longrightarrow
 \text{слабая форма}
 \longrightarrow
 \text{FEM-система},\\
 \text{FEM-система}
 &\longrightarrow
 \text{прямое состояние}
 \longrightarrow
 \text{сопряжённое состояние}
 \longrightarrow
 \nabla_\theta j.
\end{aligned}
\]
Этот путь является математическим ядром физически проверяемого генеративного проектирования. Следующие главы могут подставлять в него конкретные законы механики и теплопереноса, обучаемые суррогаты или генеративные модели, не смешивая физическую постановку, дискретизацию и оптимизацию.
